% !TEX program = xelatex
\documentclass[conference]{IEEEtran}
\IEEEoverridecommandlockouts

\usepackage{fontspec}
\setmainfont{Times New Roman}
\setsansfont{Arial}
\setmonofont[Scale=0.85]{Menlo}

\usepackage{cite}
\usepackage{amsmath,amssymb,amsfonts}
\usepackage{algorithm}
\usepackage{algorithmic}
\usepackage{graphicx}
\usepackage{booktabs}
\usepackage{array}
\usepackage{multirow}
\usepackage{xcolor}
\usepackage{hyperref}

\hypersetup{
    colorlinks=true,
    linkcolor=blue!60!black,
    citecolor=blue!60!black,
    urlcolor=blue!60!black
}

% Việt hóa nhãn chuẩn IEEE
\renewcommand{\figurename}{Hình}
\renewcommand{\tablename}{Bảng}
\renewcommand{\abstractname}{Tóm tắt}
\renewcommand{\IEEEkeywordsname}{Từ khóa}
\renewcommand{\refname}{Tài liệu tham khảo}
\floatname{algorithm}{Thuật toán}

\begin{document}

\title{Cài Đặt, Tối Ưu Hóa và Ứng Dụng Thực Tế Thuật Toán negFIN Khai Thác Tập Phổ Biến}

\author{
\IEEEauthorblockN{Nguyễn Thanh Tùng}
\IEEEauthorblockA{\textit{Khoa Công nghệ Thông tin} \\
\textit{Trường Đại học Khoa học Tự nhiên, ĐHQG-HCM}\\
TP. Hồ Chí Minh, Việt Nam \\
MSSV: 24127583}
\and
\IEEEauthorblockN{Nguyễn Tuấn Khanh}
\IEEEauthorblockA{\textit{Khoa Công nghệ Thông tin} \\
\textit{Trường Đại học Khoa học Tự nhiên, ĐHQG-HCM}\\
TP. Hồ Chí Minh, Việt Nam \\
MSSV: 24127415}
}

\maketitle

\begin{abstract}
Khai thác tập phổ biến (Frequent Itemset Mining -- FIM) là bài toán nền tảng trong lĩnh vực khai phá dữ liệu, đóng vai trò tiền đề cho sinh luật kết hợp, phân cụm và phân loại mẫu. Trong đồ án này, chúng tôi nghiên cứu, cài đặt hoàn toàn từ đầu và tối ưu hóa thuật toán \texttt{negFIN} (Aryabarzan et al., 2018) bằng ngôn ngữ \textbf{Julia} ($\ge 1.9$) dựa trên cấu trúc cây mã hóa nhị phân \texttt{BMC-tree} và tập nút phủ định \texttt{NegNodeset}. Quy trình triển khai gồm bốn bậc: phiên bản cơ bản (Baseline) bám sát giả mã gốc, bộ kiểm thử tự động (66/66 kiểm thử đạt 100\%), phiên bản tối ưu hóa chuyên biệt theo đặc trưng Julia (Multiple Dispatch bit-packing, cây LCRS, bố cục SoA, cắt sớm Early-Abort và xuất Promotion không cấp phát \texttt{FIS\_parent}), cùng giao diện dòng lệnh (CLI) đọc/ghi chuẩn SPMF. Thực nghiệm toàn diện trên 5 bộ dữ liệu chuẩn (\texttt{Chess}, \texttt{Mushroom}, \texttt{Accidents}, \texttt{Retail}, \texttt{T10I4D100K}) và dữ liệu tổng hợp cho thấy bản tối ưu đạt độ chính xác khớp 100\%, tăng tốc đến \textbf{3.26 lần} trên benchmark thực (\textbf{4.88 lần} trên giao dịch dài), giảm \textbf{27.5\%--69.5\%} bộ nhớ Heap cấp phát và mở rộng gần như tuyến tính theo số giao dịch. Cuối cùng, mã nguồn \texttt{negFIN} tự viết được tích hợp trực tiếp vào hai ứng dụng thực tế thuần Julia: (1) Phân tích giỏ hàng (\textit{Market Basket Analysis}) trên 9,835 giao dịch \texttt{Groceries} phát hiện 322 luật kết hợp với \textit{Lift} đạt tới \textbf{3.37}, và (2) Khai thác mẫu trực quan theo mô hình \textit{Bag of Visual Words} trên bộ ảnh chữ số \texttt{UCI Digits 8$\times$8} đạt độ làm giàu đặc trưng (\textit{enrichment}) tới \textbf{4.00 lần}. Toàn bộ mã nguồn, dữ liệu benchmark trên Hugging Face và notebook Google Colab đều được công khai để đảm bảo khả năng tái lập 100\%.
\end{abstract}

\begin{IEEEkeywords}
Frequent Itemset Mining, negFIN, BMC-tree, NegNodeset, Julia, Multiple Dispatch, Market Basket Analysis, Bag of Visual Words.
\end{IEEEkeywords}

% ==============================================================================
\section{Giới Thiệu}
% ==============================================================================

\subsection{Bối cảnh và Định nghĩa Bài toán}
Khai thác tập phổ biến (\textit{Frequent Itemset Mining} -- FIM), lần đầu tiên được đề xuất bởi Agrawal et al.~\cite{b1} cho bài toán phân tích giỏ hàng (\textit{Market Basket Analysis}), là một trong những tác vụ cốt lõi nhất của khai phá dữ liệu (\textit{Data Mining}). Kết quả tập phổ biến là đầu vào trực tiếp cho nhiều bài toán quan trọng khác như phát hiện luật kết hợp (\textit{Association Rule Mining}), phân cụm dữ liệu giao dịch, phân loại dựa trên mẫu phân biệt và phân tích hành vi người dùng.

Về mặt hình thức, cho tập hữu hạn các mục hàng (items) $I = \{i_1, i_2, \dots, i_{n_{it}}\}$. Một giao dịch (\textit{transaction}) $T$ là một tập con $T \subseteq I$ gắn với định danh duy nhất \texttt{TID}, và cơ sở dữ liệu giao dịch là $DB = \{T_1, T_2, \dots, T_{n_t}\}$ với kích thước $|DB| = n_t$. Một tập con $P \subseteq I$ có $|P| = k$ được gọi là một \textit{$k$-itemset}. Độ hỗ trợ (\textit{support}) của $P$ trong $DB$ được định nghĩa là:
\begin{align}
\text{sup\_count}(P) &= |\{T \in DB \mid P \subseteq T\}|, \nonumber \\
\text{support}(P) &= \frac{\text{sup\_count}(P)}{|DB|}.
\label{eq:support}
\end{align}
Với ngưỡng hỗ trợ tối thiểu do người dùng chỉ định $\text{min\_sup} \in (0, 1]$, tập $P$ được gọi là \textit{tập phổ biến} (\textit{frequent itemset}) nếu và chỉ nếu $\text{support}(P) \ge \text{min\_sup}$ (tương đương $\text{sup\_count}(P) \ge \lceil \text{min\_sup} \cdot |DB| \rceil$). Bài toán FIM yêu cầu tìm toàn bộ tập phổ biến $\mathcal{F} = \{(P, \text{sup\_count}(P)) \mid P \subseteq I, P \neq \emptyset, \text{support}(P) \ge \text{min\_sup}\}$. Do không gian tìm kiếm có kích thước $2^{|I|} - 1$, bài toán FIM thuộc lớp khó tính toán (NP-hard) khi $|I|$ lớn hoặc $\text{min\_sup}$ nhỏ.

\subsection{Tổng quan các Hướng tiếp cận và Động lực chọn \texttt{negFIN}}
Các thuật toán FIM kinh điển được chia thành ba nhánh lớn:
\begin{enumerate}
    \item \textbf{Sinh ứng viên (\textit{Candidate Generation}):} Tiêu biểu là \texttt{Apriori}~\cite{b2}, sử dụng tính chất phản đơn điệu (\textit{anti-monotone}) để cắt tỉa ứng viên nhưng phải quét CSDL nhiều lần, kém hiệu quả khi mẫu dài.
    \item \textbf{Tăng trưởng mẫu (\textit{Pattern Growth}) và Khai thác dọc (\textit{Vertical Mining}):} \texttt{FP-Growth}~\cite{b3} nén CSDL vào cây \texttt{FP-tree} chỉ với 2 lần quét; \texttt{Eclat}~\cite{b4} biểu diễn giao dịch theo dạng dọc (\texttt{tidset}/\texttt{diffset}). Tuy nhiên, \texttt{FP-Growth} tốn chi phí xây cây điều kiện trên dữ liệu thưa, còn \texttt{Eclat} gặp chi phí giao tập \texttt{tidset} lớn trên dữ liệu dày đặc.
    \item \textbf{Họ cấu trúc tập nút trên cây tiền tố (\textit{Prefix-tree Nodeset}):} Mở đầu với \texttt{Node-list} (PPV)~\cite{b5} và \texttt{N-list} (PrePost)~\cite{b6} lưu cả thứ tự duyệt tiền tố (\textit{pre-order}) lẫn hậu tố (\textit{post-order}) trên cây PPC-tree; tiếp theo là \texttt{Nodeset} (FIN)~\cite{b7} giảm một nửa bộ nhớ khi chỉ lưu một thứ tự; và \texttt{DiffNodeset} (dFIN)~\cite{b8} dùng phép hiệu giữa hai tập nút để thu nhỏ kích thước danh sách nút.
\end{enumerate}

Dù \texttt{dFIN} nhanh hơn các thuật toán trước đó, phép lấy hiệu giữa hai \texttt{DiffNodeset} có độ phức tạp $O(m + n)$ (với $m, n$ là kích thước hai tập nút) vẫn là nút cổ chai lớn. Năm 2018, Aryabarzan et al.~\cite{b9} đề xuất cấu trúc \textbf{\texttt{NegNodeset}} dựa trên cây mã hóa bit \textbf{\texttt{BMC-tree}} và thuật toán \textbf{\texttt{negFIN}}, rút gọn độ phức tạp trích xuất mỗi tập nút mới xuống chỉ còn $O(n)$ nhờ các phép toán bitwise cấp độ phần cứng CPU ($O(1)$) kết hợp chiến lược cắt tỉa \textit{Promotion}.

\subsection{Đóng góp của Đồ án}
Đồ án thực hiện trọn vẹn ba mục tiêu trọng tâm:
\begin{itemize}
    \item \textbf{Chương 1:} Cài đặt hoàn chỉnh thuật toán \texttt{negFIN} theo 4 bậc (Baseline, Unit Tests tự động, bản tối ưu hóa Julia-native và CLI chuẩn SPMF), chứng minh hiệu quả tối ưu hóa định lượng.
    \item \textbf{Chương 2:} Thực thi 6 thực nghiệm chuẩn khoa học trên 5 bộ dữ liệu benchmark kinh điển từ kho FIMI/SPMF (\texttt{Chess}, \texttt{Mushroom}, \texttt{Accidents}, \texttt{Retail}, \texttt{T10I4D100K}) được nhóm lưu trữ công khai trên Hugging Face nhằm đảm bảo khả năng tái lập, cùng dữ liệu tổng hợp, phân tích sâu theo lý thuyết bài báo gốc và đề xuất 3 hướng cải tiến tiếp theo.
    \item \textbf{Chương 3:} Tích hợp trực tiếp mã nguồn \texttt{negFIN} tự viết vào hai ứng dụng thực tế(\textit{Market Basket Analysis} trên \texttt{Groceries} và \textit{Bag of Visual Words} trên ảnh \texttt{UCI Digits 8$\times$8}), đảm bảo tái lập 100\% qua Google Colab.
\end{itemize}

% ==============================================================================
\section{Chương 1: Cơ Sở Lý Thuyết và Cài Đặt Thuật Toán \texttt{negFIN} Bằng Julia}
% ==============================================================================

\subsection{Cơ sở Lý thuyết Thuật toán \texttt{negFIN}}
Dựa trên bài báo gốc của Aryabarzan et al.~\cite{b9}, thuật toán \texttt{negFIN} xoay quanh bốn khái niệm cốt lõi sau:

\subsubsection{Thứ tự mục hàng $L_1$ và Mã hóa Bitmap (\texttt{BMC})}
Sau lần quét CSDL đầu tiên, ta thu được tập các 1-itemset phổ biến $F_1$ với $n_f = |F_1|$. Định nghĩa quan hệ thứ tự $\succ$ trên $F_1$: $i_2 \succ i_1$ nếu và chỉ nếu $\text{support}(i_2) \ge \text{support}(i_1)$ (phá hòa nhất quán theo mã item). Vector $L_1 = [i_0, i_1, \dots, i_{n_f - 1}]$ sắp xếp các phần tử của $F_1$ theo thứ tự \textbf{support không giảm} ($i_{n_f - 1} \succ \dots \succ i_1 \succ i_0$). Mỗi item $i \in L_1$ được gán chỉ số bit 0-based $\text{index}(i) \in \{0, \dots, n_f - 1\}$.

Mỗi itemset $P_k = i_k i_{k-1} \dots i_1$ ($i_k \succ \dots \succ i_1$) được mã hóa bởi chuỗi bit độ dài $n_f$:
\begin{align}
\text{BMC}(P_k) &= b_{n_f - 1} b_{n_f - 2} \dots b_1 b_0, \nonumber \\
b_j &=
\begin{cases}
1 & \text{nếu } L_1[j] \in P_k, \\
0 & \text{ngược lại.}
\end{cases}
\label{eq:bmc}
\end{align}

\subsubsection{Cây Mã hóa Bitmap (\texttt{BMC-tree})}
\texttt{BMC-tree} là cây tiền tố (\textit{prefix tree}) có nút gốc rỗng $\emptyset$ với $\text{root.bitmap\_code} = 0$ (\textit{Property 3}). Với mỗi giao dịch $T \in DB$, các item không phổ biến bị loại bỏ và các item còn lại được sắp xếp theo thứ tự \textbf{ngược của $L_1$} (support giảm dần) trước khi chèn vào cây. Mỗi nút $N$ chứa 4 trường: \texttt{item\_name}, \texttt{count}, \texttt{bitmap\_code} và danh sách con. Theo \textit{Property 4}, mã bit của nút con $N$ chứa item $i$ được tính trực tiếp từ nút cha bằng một phép toán \texttt{OR} nhị phân:
\begin{equation}
N.\text{bitmap\_code} = N.\text{father.bitmap\_code} \lor 2^{\text{index}(i)}.
\label{eq:prop4}
\end{equation}
Nếu $m = \text{index}(i)$, đoạn bit $b_{n_f - 1} \dots b_m$ gọi là \textit{phần chính (main section)} -- trong đó bit $0$ khẳng định chắc chắn item tương ứng không nằm trên đường đi từ gốc tới $N$ (\textit{Property 2}) -- còn đoạn $b_{m-1} \dots b_0$ gọi là \textit{phần không quan tâm (don't-care section)} và luôn bằng $0$ (\textit{Property 1}). Sau khi duyệt cây \texttt{BMC-tree} một lần để thu thập cặp thông tin $N\text{-info} = (\text{bitmap\_code}, \text{count})$ vào $\text{Nodeset}(i)$ cho từng $i \in L_1$, toàn bộ cây \texttt{BMC-tree} được giải phóng.

\subsubsection{Cấu trúc \texttt{NegNodeset} và Tính Support trong $O(n)$}
Với một $k$-itemset $P_k = i_k P_{k-1}$ ($k \ge 2$), ký hiệu $P'_k = \neg i_k P_{k-1}$ biểu thị các giao dịch chứa $P_{k-1}$ nhưng \textit{không} chứa $i_k$. Tập $\text{NegNodeset}(P_k)$ được định nghĩa là $\text{Nodeset}(P'_k)$:
\begin{itemize}
    \item \textbf{Tầng $k = 2$ ($R = i P$, $P = \{j\}$, $i \succ j$):} Trích từ $\text{Nodeset}(j)$ các nút $ni$ có bit ứng với $i$ bằng $0$ (\textit{Definition 10}):
    \begin{align}
    \text{NegNodeset}(ij) = \{&ni \in \text{Nodeset}(j) \mid \nonumber \\
    &(ni.\text{bitmap} \land 2^{\text{index}(i)}) = 0\}.
    \label{eq:level2}
    \end{align}
    \item \textbf{Tầng $k \ge 3$ ($R = i_k i_{k-1} P_{k-2}$, đặt $Q = i_k P_{k-2}$):} Trích trực tiếp từ một tập duy nhất $\text{NegNodeset}(Q)$ (nút anh em bên phải ở cùng mức cha) các nút $ni$ có bit ứng với $i_{k-1}$ bằng $1$ (\textit{Property 7}):
    \begin{align}
    \text{NegNodeset}(R) = \{&ni \in \text{NegNodeset}(Q) \mid \nonumber \\
    &(ni.\text{bitmap} \land 2^{\text{index}(i_{k-1})}) \neq 0\}.
    \label{eq:levelk}
    \end{align}
\end{itemize}
Từ \textit{Property 6} và \textit{Property 8}, độ hỗ trợ của $R$ được tính chính xác qua hiệu:
\begin{equation}
\text{sup\_count}(R) = \text{sup\_count}(P) - \sum_{ni \in \text{NegNodeset}(R)} ni.\text{count}.
\label{eq:sup_diff}
\end{equation}

\subsubsection{Cắt tỉa Tương đương (\textit{Superset Equivalence / Promotion})}
Theo \textit{Property 9}, nếu $\text{sup\_count}(R) = \text{sup\_count}(P)$ (tương đương $\text{NegNodeset}(R) = \emptyset$), thì với mọi tập $Q$ rời nhau với $P$ và $i \notin Q$, ta luôn có $\text{sup\_count}(P \cup Q \cup \{i\}) = \text{sup\_count}(P \cup Q)$. Khi đó, item $i$ được đưa vào tập \texttt{equivalent\_items} của nút hiện tại mà \textbf{không} tạo nút con trên cây liệt kê (\textit{Set-Enumeration Tree}). Mọi tập con của \texttt{equivalent\_items} kết hợp với đường đi hiện tại đều có chung support và được sinh trực tiếp, giúp cắt giảm đáng kể không gian tìm kiếm trên dữ liệu dày đặc.

\subsubsection{Phân tích Độ phức tạp và So sánh họ Thuật toán Nodeset}
Gọi $l$ là tổng số nút ở các tầng $k \ge 2$ của cây liệt kê tập phổ biến và $n$ là kích thước trung bình của tập nút nguồn. Vì mỗi lần trích \texttt{NegNodeset} chỉ quét đúng một tập nút nguồn với phép kiểm tra bit $O(1)$, độ phức tạp pha khai thác của \texttt{negFIN} là $O(l \cdot n)$, giảm rõ rệt so với $O(l(m+n))$ của \texttt{dFIN}~\cite{b8} (trong đó $n \le m$). Bảng~\ref{tab:nodeset_family} đối chiếu sự tiến hóa của 5 thế hệ thuật toán dựa trên cây tiền tố.

\begin{table}[htbp]
\caption{So sánh 5 thế hệ thuật toán khai thác tập phổ biến dựa trên cấu trúc cây tiền tố}
\label{tab:nodeset_family}
\centering
\resizebox{\columnwidth}{!}{%
\begin{tabular}{@{}lllll@{}}
\toprule
\textbf{Thuật toán} & \textbf{Cấu trúc} & \textbf{Mã hóa nút} & \textbf{Phép toán} & \textbf{Độ phức tạp} \\
\midrule
PPV~\cite{b5} & Node-list & Pre + Post & Giao 2 tập & $O(l(m+n))$ \\
PrePost~\cite{b6} & N-list & Pre + Post & Giao 2 tập & $O(l(m+n))$ \\
FIN~\cite{b7} & Nodeset & Pre (hoặc Post) & Giao 2 tập & $O(l(m+n))$ \\
dFIN~\cite{b8} & DiffNodeset & Pre (hoặc Post) & Hiệu 2 tập & $O(l(m+n))$ \\
\textbf{negFIN}~\cite{b9} & \textbf{NegNodeset} & \textbf{Bitmap (BMC)} & \textbf{Bitwise 1 tập} & \textbf{$O(l \cdot n)$} \\
\bottomrule
\end{tabular}%
}
\end{table}

\subsubsection{Ví dụ Minh họa Từng bước (Table 1 bài báo gốc)}
Xét CSDL mẫu gồm 5 giao dịch trong Bảng~\ref{tab:example1} với ngưỡng $\text{min\_sup} = 0.4$ (tức $\text{min\_sup\_count} = \lceil 0.4 \times 5 \rceil = 2$). Sau lần quét đầu tiên, 5 item phổ biến được sắp xếp tăng dần theo support thành $L_1 = [e(0), d(1), c(2), b(3), a(4)]$ (chỉ số trong ngoặc là vị trí bit 0-based).

\begin{table}[htbp]
\caption{CSDL giao dịch mẫu (Table 1, Aryabarzan et al., 2018) với $\text{min\_sup}=0.4$}
\label{tab:example1}
\centering
\begin{tabular}{@{}cll@{}}
\toprule
\textbf{TID} & \textbf{Các item gốc} & \textbf{Item phổ biến sắp theo $L_1^{\text{reverse}}$} \\
\midrule
1 & \texttt{e, b, g, d} & \texttt{b, d, e} \\
2 & \texttt{c, e, b, a} & \texttt{a, b, c, e} \\
3 & \texttt{c, b, a, i} & \texttt{a, b, c} \\
4 & \texttt{a, d, h} & \texttt{a, d} \\
5 & \texttt{a, d, c, b, f} & \texttt{a, b, c, d} \\
\bottomrule
\end{tabular}
\end{table}

Khi chèn các giao dịch đã sắp theo $L_1^{\text{reverse}}$ vào \texttt{BMC-tree} và trích \texttt{Nodeset} cho item $d$ (bit 1), ta thu được:
\begin{align*}
\text{Nodeset}(d) &= \{(\texttt{01010}, 1), (\texttt{11110}, 1), (\texttt{10010}, 1)\}, \\
\text{sup}(d) &= 1 + 1 + 1 = 3.
\end{align*}
Ở tầng 2 (\textit{Definition 10}), lọc các nút có bit $a$ (bit 4) bằng $0$ cho $\text{NegNodeset}(ad) = \{(\texttt{01010}, 1)\} \Rightarrow \text{sup}(ad) = 3 - 1 = 2$. Tương tự với bit $b$ (bit 3) bằng $0$, ta có $\text{NegNodeset}(bd) = \{(\texttt{10010}, 1)\} \Rightarrow \text{sup}(bd) = 3 - 1 = 2$. Ở tầng 3 (\textit{Property 7}), để tính $R = abd$ từ nút $P = bd$ và anh em bên phải $Q = ad$, ta chỉ cần kiểm tra bit $b$ (bit 3) bằng $1$ trên $\text{NegNodeset}(ad)$: nút $(\texttt{01010}, 1)$ có bit $b = 1$ nên $\text{NegNodeset}(abd) = \{(\texttt{01010}, 1)\}$, suy ra $\text{sup}(abd) = \text{sup}(bd) - 1 = 1 < 2$ (bị cắt tỉa). Đối với nút $c$, mọi nút thuộc $\text{Nodeset}(c) = \{(\texttt{11100}, 3)\}$ đều có bit $a=1$ và $b=1$ nên $\text{NegNodeset}(ac) = \text{NegNodeset}(bc) = \emptyset$; cả $a$ và $b$ được đưa vào \texttt{equivalent\_items} của $c$ (\textit{Promotion}), sinh trực tiếp 4 tập $\{c\}, \{ac\}, \{bc\}, \{abc\}$ cùng có $\text{sup} = 3$ mà không cần đệ quy thêm. Tổng cộng thu được đúng 12 tập phổ biến khớp 100\% với kết quả vét cạn và CLI \texttt{src/main.jl}.

\begin{algorithm}[t]
\caption{Khung thuật toán \texttt{negFIN} (phiên bản tối ưu)}
\label{alg:negfin}
\begin{algorithmic}[1]
\REQUIRE Cơ sở dữ liệu $DB$, ngưỡng $\text{min\_sup}$
\ENSURE Tập tất cả tập phổ biến $\mathcal{F}$
\STATE Quét $DB$ tìm $F_1$, sắp xếp $L_1$ tăng dần theo support, $n_f \leftarrow |L_1|$
\STATE Chọn kiểu bitmap $B \in \{\texttt{UInt64}, \texttt{UInt128}, \texttt{NTuple}, \texttt{Vector\{UInt64\}}\}$ theo $n_f$
\STATE Xây cây \texttt{BMC-tree} dạng LCRS với mã bit kiểu $B$; trích $\text{Nodeset}[1..n_f]$ dạng SoA; giải phóng cây
\FOR{$p = 1$ \TO $n_f$}
    \STATE $\text{max\_neg} \leftarrow \text{sup}[p] - \text{min\_sup\_count}$; khởi tạo \texttt{children}, \texttt{equiv\_stack}
    \FOR{$i = p + 1$ \TO $n_f$}
        \STATE Quét $\text{Nodeset}[p]$ kiểm tra bit $i = 0$, cộng dồn $\text{neg\_sup}$
        \IF{$\text{neg\_sup} > \text{max\_neg}$}
            \STATE \textbf{break} \COMMENT{Cắt sớm Early-Abort}
        \ENDIF
        \IF{$\text{neg\_sup} = 0$}
            \STATE Đẩy $L_1[i]$ vào \texttt{equiv\_stack} \COMMENT{Promotion (Property 9)}
        \ELSE
            \STATE Thêm nút con $(i, \text{sup}[p] - \text{neg\_sup}, \text{NegNodeset})$
        \ENDIF
    \ENDFOR
    \STATE Xuất đệ quy mọi tổ hợp $\text{prefix} \cup S$ ($S \subseteq \text{equiv\_stack}$)
    \STATE Gọi đệ quy mở rộng các nút con theo anh em bên phải (\textit{Property 7})
    \STATE Giải phóng $\text{Nodeset}[p]$
\ENDFOR
\end{algorithmic}
\end{algorithm}

\subsection{Kiến trúc Cài đặt 4 Bậc và Tối ưu hóa Đặc trưng Julia}
Toàn bộ mã nguồn Chương 1 được tổ chức chuẩn hóa trong thư mục \texttt{src/} tuân thủ chặt chẽ quy chuẩn Julia BlueStyle (định kiểu tường minh, thụt lề 4 khoảng trắng, $\le 92$ ký tự/dòng, không dùng biến toàn cục không \texttt{const}):
\begin{itemize}
    \item \texttt{src/structures.jl}: Định nghĩa cấu trúc dữ liệu cho cả bản cơ bản và bản tối ưu.
    \item \texttt{src/algorithm/negfin\_baseline.jl}: \textbf{Bậc 1 (Baseline)} -- bám sát giả mã Algorithm 1--4 của bài báo gốc, sử dụng \texttt{BitVector}, danh sách con \texttt{Vector\{BMCNodeBaseline\}} và tập \texttt{FIS\_parent} tường minh.
    \item \texttt{tests/runtests.jl}, \texttt{tests/test\_correctness.jl}, \texttt{tests/test\_benchmark.jl}: \textbf{Bậc 2 (Unit Tests)} -- kiểm thử tự động 100\% trên $\ge 7$ kịch bản khác nhau (ví dụ Table 1 trong bài báo gốc, CSDL rỗng, $\text{min\_sup}$ quá cao, Promotion cực đại $2^8 - 1 = 255$ tập, giao dịch và item trùng lặp, mã item không liên tục rất lớn, kiểm tra đủ 5 nhánh Multiple Dispatch và vòng khép kín I/O SPMF), kèm kiểm chứng độc lập bằng thuật toán vét cạn và cầu nối đối chiếu \texttt{spmf.jar}.
    \item \texttt{src/algorithm/negfin\_optimized.jl}: \textbf{Bậc 3 (Optimized)} -- áp dụng 5 kỹ thuật tối ưu hóa đặc trưng của Julia:
    \begin{enumerate}
        \item \textbf{Multiple Dispatch \& Function Barrier theo $n_f$:} Chuyên biệt hóa kiểu bitmap $B$ tại thời điểm biên dịch: \texttt{UInt64} ($n_f \le 64$), \texttt{UInt128} ($n_f \le 128$, vừa khít \texttt{Chess} và \texttt{Mushroom} ngay trên thanh ghi CPU), \texttt{NTuple\{3,UInt64\}} ($n_f \le 192$), \texttt{NTuple\{4,UInt64\}} ($n_f \le 256$), và \texttt{Vector\{UInt64\}} ($n_f > 256$, như \texttt{Retail}, \texttt{Accidents}, \texttt{T10I4D100K}). Loại bỏ hoàn toàn dynamic dispatch và kiểu trừu tượng \texttt{Any} trong vòng lặp nóng.
        \item \textbf{Cây \texttt{BMC-tree} dạng Left-Child Right-Sibling (LCRS):} Cấu trúc \texttt{BMCNodeOpt\{B\}} dùng hai con trỏ \texttt{first\_child} và \texttt{next\_sibling} thay vì mảng \texttt{Vector} con, loại bỏ hàng trăm nghìn lần cấp phát mảng rỗng cho các nút lá trên dữ liệu thưa.
        \item \textbf{Bố cục Structure-of-Arrays (SoA) và \texttt{@inbounds}:} Tách \texttt{NodeSetOpt\{B\}} thành hai mảng phẳng liên tiếp \texttt{bitmaps::Vector\{B\}} và \texttt{counts::Vector\{Int32\}}, kết hợp tính sẵn mặt nạ bit (\texttt{bit\_locator}) ngoài vòng lặp nóng giúp tối ưu bộ nhớ đệm CPU (cache locality).
        \item \textbf{Cắt sớm (\textit{Early-Abort Pruning}):} Trong vòng lặp lọc bit trích \texttt{NegNodeset}, cộng dồn ngay \texttt{neg\_sup}; nếu $\text{neg\_sup} > \text{parent\_sup} - \text{min\_sup\_count}$ thì thoát vòng lặp lập tức (dòng 8--9, Thuật toán~\ref{alg:negfin}).
        \item \textbf{Xuất Promotion không cấp phát \texttt{FIS\_parent}:} Thay vì lưu tập hợp \texttt{FIS\_parent} có thể phình to theo hàm mũ, sử dụng hai ngăn xếp phẳng \texttt{prefix} và \texttt{equiv\_stack} cùng hàm callback streaming \texttt{emit\_fn}, đồng thời giải phóng sớm \texttt{nodesets[pos]} ngay sau khi duyệt xong từng nhánh tầng 1.
    \end{enumerate}
    \item \texttt{src/utils.jl} và \texttt{src/main.jl}: \textbf{Bậc 4 (I/O \& CLI)} -- đọc file \texttt{.txt} chuẩn SPMF (tự bỏ qua dòng trống/metadata và khử item trùng trong cùng giao dịch), xuất đúng định dạng \texttt{item1 item2 \#SUP: count}, hỗ trợ tham số dòng lệnh \texttt{--input}, \texttt{--output}, \texttt{--minsup}, \texttt{--algo baseline|optimized} và \texttt{--stats}.
\end{itemize}

\subsection{Kết quả Kiểm thử Tự động và Đối chứng Sơ bộ}
Toàn bộ hệ thống kiểm thử đạt kết quả \textbf{56/56 PASS (100\%)}: gồm \textbf{29/29 PASS} ở bộ kiểm thử Chương 1--2 (\texttt{tests/runtests.jl}, thời gian 4.1 giây) và \textbf{27/27 PASS} ở bộ kiểm thử tích hợp Chương 3 (\texttt{tests/test\_ch3\_julia.jl}, thời gian 5.2 giây). Bảng~\ref{tab:unit_tests} liệt kê chi tiết các kịch bản kiểm thử bao phủ đầy đủ trường hợp biên và đa nhánh dispatch, còn Bảng~\ref{tab:ch1_bench} đối chứng trực tiếp tốc độ và bộ nhớ cấp phát giữa \texttt{negfin\_baseline} và \texttt{negfin\_optimized} sau bước warm-up JIT.

\begin{table}[htbp]
\caption{Tổng hợp bộ kiểm thử tự động Chương 1--3 (56/56 PASS)}
\label{tab:unit_tests}
\centering
\resizebox{\columnwidth}{!}{%
\begin{tabular}{@{}llc@{}}
\toprule
\textbf{Nhóm kiểm thử} & \textbf{Kịch bản / Trường hợp biên} & \textbf{Kết quả} \\
\midrule
1. Paper Example 1 & Table 1 (5 tx, $\text{min\_sup}=0.4$), 12 itemsets & 4/4 PASS \\
2. Edge Cases & CSDL rỗng; $\text{min\_sup}$ quá cao (0 tập) & 6/6 PASS \\
3. Extreme Promotion & 8 item đồng xuất hiện ($2^8-1=255$ tập) & 3/3 PASS \\
4. Duplicates & Giao dịch trùng \& item trùng trong dòng & 2/2 PASS \\
5. Sparse IDs & Mã item không liên tục, rất lớn ($10^6$) & 2/2 PASS \\
6. Multi-Width Dispatch & $n_f \in \{40, 80, 150, 220, 300\}$ (5 kiểu $B$) & 5/5 PASS \\
7. SPMF I/O Round-Trip & Ghi chuẩn \texttt{\#SUP:} và đọc đối chiếu lại & 1/1 PASS \\
8. Benchmark Suite & Đối chứng Baseline vs. Optimized (3 tập) & 6/6 PASS \\
9. Chapter 3 Integration & Market Basket, Rules, K-means, Visual BoVW & 27/27 PASS \\
\midrule
\textbf{Tổng cộng} & \textbf{Toàn bộ mã nguồn Julia} & \textbf{56/56 PASS} \\
\bottomrule
\end{tabular}%
}
\end{table}

\begin{table}[htbp]
\caption{Đối chứng tốc độ và bộ nhớ cấp phát giữa Baseline và Optimized (Unit Benchmark sau warm-up JIT)}
\label{tab:ch1_bench}
\centering
\resizebox{\columnwidth}{!}{%
\begin{tabular}{@{}lrrrrrr@{}}
\toprule
\textbf{Tập kiểm thử (Kiểu $B$)} & \textbf{Base} & \textbf{Opt} & \textbf{Tăng} & \textbf{Base} & \textbf{Opt} & \textbf{Giảm} \\
 & \textbf{(ms)} & \textbf{(ms)} & \textbf{tốc} & \textbf{(MB)} & \textbf{(MB)} & \textbf{Heap} \\
\midrule
Dense-64 (\texttt{UInt64}) & 15.96 & 3.06 & \textbf{5.21$\times$} & 33.61 & 11.35 & -66.2\% \\
Dense-120 (\texttt{UInt128}) & 25.33 & 4.84 & \textbf{5.23$\times$} & 86.00 & 24.89 & -71.1\% \\
Sparse-400 (\texttt{Vec\{UInt64\}}) & 3.63 & 2.27 & \textbf{1.60$\times$} & 13.33 & 3.39 & -74.6\% \\
\bottomrule
\end{tabular}%
}
\end{table}

% ==============================================================================
\section{Chương 2: Thực Nghiệm và Đánh Giá Hiệu Năng}
% ==============================================================================

\subsection{Môi trường, Quy trình Đo và Tập dữ liệu Benchmark}
Toàn bộ thực nghiệm Chương 2 được thực thi tự động qua \texttt{scripts/run\_benchmarks.jl} và notebook thuần Julia \texttt{ch2\_experiments.ipynb} (có hai phiên bản \texttt{notebooks/colab/} và \texttt{notebooks/local/}) trên máy Apple Silicon (macOS, Julia 1.13.3). Để đảm bảo tính khoa học và tái lập:
\begin{itemize}
    \item Mọi phép đo thời gian đều thực hiện \textbf{1 lần chạy warm-up JIT} trước khi bấm giờ, sau đó lấy \textbf{trung vị (median) của 3 lần chạy độc lập} kèm \texttt{GC.gc()} giữa các lần.
    \item Bộ nhớ được đo đồng thời bằng \texttt{@allocated} (tổng Heap cấp phát) và \textbf{Peak RSS} tiến trình con qua \texttt{/usr/bin/time -l}.
    \item Mọi bước sinh dữ liệu tổng hợp đều cố định hạt giống \texttt{MersenneTwister(42 + i)}.
\end{itemize}

Chúng tôi đánh giá trên đầy đủ \textbf{5/5 bộ dữ liệu benchmark chuẩn} có nguồn gốc từ kho FIMI/SPMF (trùng khớp Table 2 của bài báo gốc~\cite{b9}), được giữ nguyên định dạng SPMF và lưu trữ công khai (mirror) trên Hugging Face (\url{https://huggingface.co/datasets/stephent9/Lab1_Frequent_Itemset_Mining}) để phục vụ tải tự động trong quá trình tái thực hiện, tóm tắt tại Bảng~\ref{tab:datasets}.

\begin{table}[htbp]
\caption{Đặc trưng 5 bộ dữ liệu Benchmark chuẩn SPMF sử dụng trong Chương 2}
\label{tab:datasets}
\centering
\resizebox{\columnwidth}{!}{%
\begin{tabular}{@{}lcrrrr@{}}
\toprule
\textbf{Tập dữ liệu} & \textbf{Loại} & \textbf{\#Giao dịch} & \textbf{\#Items} & \textbf{Dài TB} & \textbf{Dung lượng} \\
\midrule
\texttt{Chess} & Dense & 3,196 & 75 & 37.00 & 334 KB \\
\texttt{Mushroom} & Dense & 8,416 & 119 & 23.00 & 589 KB \\
\texttt{Accidents} & Dense & 340,183 & 468 & 33.81 & 34.0 MB \\
\texttt{Retail} & Sparse & 88,162 & 16,469 & 10.31 & 4.0 MB \\
\texttt{T10I4D100K} & Sparse & 100,000 & 949 & 10.10 & 3.8 MB \\
\bottomrule
\end{tabular}%
}
\end{table}

\subsection{Thực nghiệm 1: Kiểm tra Tính Đúng Đắn (\textit{Correctness})}
Bảng~\ref{tab:exp1} tổng hợp kết quả kiểm chứng trên cả 5 bộ dữ liệu chuẩn. Toàn bộ số lượng tập phổ biến và giá trị \texttt{support} từng tập giữa \texttt{negfin\_baseline} và \texttt{negfin\_optimized} đều \textbf{khớp tuyệt đối 100\%} (không có tập trùng lặp hay sai lệch).

\begin{table}[htbp]
\caption{Kết quả Thực nghiệm 1: Kiểm tra tính đúng đắn trên 5 bộ Benchmark}
\label{tab:exp1}
\centering
\resizebox{\columnwidth}{!}{%
\begin{tabular}{@{}lrrrrl@{}}
\toprule
\textbf{Tập dữ liệu} & \textbf{min\_sup} & \textbf{\#Tx} & \textbf{Base} & \textbf{Opt} & \textbf{Tỷ lệ khớp} \\
\midrule
\texttt{Chess} & 0.80 & 3,196 & 8,227 & 8,227 & \textbf{MATCH 100\%} \\
\texttt{Mushroom} & 0.30 & 8,416 & 2,587 & 2,587 & \textbf{MATCH 100\%} \\
\texttt{Accidents} & 0.70 & 340,183 & 529 & 529 & \textbf{MATCH 100\%} \\
\texttt{Retail} & 0.02 & 88,162 & 55 & 55 & \textbf{MATCH 100\%} \\
\texttt{T10I4D100K} & 0.02 & 100,000 & 155 & 155 & \textbf{MATCH 100\%} \\
\bottomrule
\end{tabular}%
}
\end{table}

\subsection{Thực nghiệm 2: Thời gian Chạy theo \texttt{min\_sup}}
Hình~\ref{fig:runtime} và Bảng~\ref{tab:exp2_full} thể hiện thời gian chạy trên 6 mốc \texttt{min\_sup} giảm dần cho từng bộ dữ liệu. Bản tối ưu luôn nhanh hơn bản cơ bản trên toàn bộ 30 điểm đo:
\begin{itemize}
    \item Trên \texttt{T10I4D100K}, \texttt{negfin\_optimized} đạt tốc độ vượt trội nhất: tăng tốc \textbf{3.00$\times$} ở $\text{min\_sup}=2\%$ (0.413s $\rightarrow$ 0.138s) và \textbf{3.26$\times$} ở $\text{min\_sup}=1.5\%$ (0.765s $\rightarrow$ 0.235s), nhờ cơ chế \textit{Early-Abort} loại bỏ sớm các ứng viên không phổ biến trên cây tìm kiếm rộng.
    \item Trên \texttt{Chess} và \texttt{Mushroom}, nhờ $n_f \le 128$ khớp hoàn toàn vào kiểu \texttt{UInt128} cấp phát trên stack, bản tối ưu xử lý 111,239 tập phổ biến (\texttt{Chess} 65\%) chỉ trong \textbf{9.8 ms} (tăng tốc 1.51$\times$--1.72$\times$) và 99,071 tập phổ biến (\texttt{Mushroom} 15\%) trong \textbf{8.3 ms} (tăng tốc 1.56$\times$).
    \item Trên tập lớn \texttt{Accidents} (340,183 giao dịch), bản tối ưu chỉ mất \textbf{0.303s--0.413s} cho cả 6 mốc \texttt{min\_sup}.
\end{itemize}

\begin{figure*}[t]
\centering
\includegraphics[width=0.96\textwidth]{figures/fig1_runtime_vs_minsup.png}
\caption{Thực nghiệm 2: So sánh thời gian chạy theo ngưỡng \texttt{min\_sup} giữa \texttt{negFIN Baseline} và \texttt{negFIN Optimized} trên 5 bộ dữ liệu benchmark (trục hoành giảm dần từ trái sang phải).}
\label{fig:runtime}
\end{figure*}

\begin{table*}[htbp]
\caption{Kết quả chi tiết Thực nghiệm 2 \& 3: Thời gian chạy (trung vị 3 lần), số tập phổ biến và bộ nhớ cấp phát trên cả 5 bộ Benchmark}
\label{tab:exp2_full}
\centering
\setlength{\tabcolsep}{5pt}
\begin{tabular}{@{}lcrrrrrrr@{}}
\toprule
\textbf{Tập dữ liệu} & \textbf{min\_sup} & \textbf{\#Itemsets} & \textbf{Base (s)} & \textbf{Opt (s)} & \textbf{Speedup} & \textbf{Alloc B (MB)} & \textbf{Alloc O (MB)} & \textbf{Giảm Heap} \\
\midrule
\multirow{6}{*}{\texttt{Chess} (Dense)}
 & 90.0\% & 622 & 0.0026 & 0.0024 & 1.10$\times$ & 0.47 & 0.27 & -41.4\% \\
 & 85.0\% & 2,669 & 0.0033 & 0.0029 & 1.14$\times$ & 1.78 & 1.10 & -38.2\% \\
 & 80.0\% & 8,227 & 0.0056 & 0.0038 & 1.48$\times$ & 4.85 & 3.04 & -37.3\% \\
 & 75.0\% & 20,993 & 0.0069 & 0.0040 & \textbf{1.72$\times$} & 11.75 & 7.42 & -36.9\% \\
 & 70.0\% & 48,731 & 0.0104 & 0.0069 & 1.51$\times$ & 25.26 & 16.62 & -34.2\% \\
 & 65.0\% & 111,239 & 0.0149 & 0.0098 & 1.51$\times$ & 55.46 & 36.83 & -33.6\% \\
\midrule
\multirow{6}{*}{\texttt{Mushroom} (Dense)}
 & 40.0\% & 505 & 0.0048 & 0.0039 & 1.24$\times$ & 0.45 & 0.20 & -55.4\% \\
 & 35.0\% & 1,121 & 0.0052 & 0.0038 & 1.36$\times$ & 1.13 & 0.45 & -60.3\% \\
 & 30.0\% & 2,587 & 0.0053 & 0.0040 & 1.33$\times$ & 1.93 & 0.88 & -54.5\% \\
 & 25.0\% & 5,393 & 0.0063 & 0.0044 & 1.44$\times$ & 4.70 & 1.88 & -60.0\% \\
 & 20.0\% & 53,337 & 0.0097 & 0.0064 & 1.51$\times$ & 18.15 & 10.11 & -44.3\% \\
 & 15.0\% & 99,071 & 0.0130 & 0.0083 & \textbf{1.56$\times$} & 37.03 & 20.71 & -44.1\% \\
\midrule
\multirow{6}{*}{\texttt{Accidents} (Dense)}
 & 75.0\% & 313 & 0.3417 & 0.3033 & 1.13$\times$ & 1.45 & 0.53 & -63.3\% \\
 & 70.0\% & 529 & 0.3585 & 0.3044 & \textbf{1.18$\times$} & 2.97 & 1.00 & -66.2\% \\
 & 65.0\% & 1,093 & 0.3735 & 0.3202 & 1.17$\times$ & 7.78 & 2.72 & -65.0\% \\
 & 60.0\% & 2,074 & 0.4101 & 0.3800 & 1.08$\times$ & 25.61 & 7.80 & -69.5\% \\
 & 55.0\% & 4,051 & 0.4241 & 0.3774 & 1.12$\times$ & 47.59 & 15.85 & -66.7\% \\
 & 50.0\% & 8,057 & 0.4831 & 0.4135 & 1.17$\times$ & 103.33 & 35.04 & -66.1\% \\
\midrule
\multirow{6}{*}{\texttt{Retail} (Sparse)}
 & 5.0\% & 16 & 0.0462 & 0.0343 & 1.35$\times$ & 1.45 & 1.10 & -24.1\% \\
 & 3.0\% & 32 & 0.0475 & 0.0362 & 1.31$\times$ & 1.81 & 1.20 & -33.4\% \\
 & 2.0\% & 55 & 0.0522 & 0.0423 & 1.23$\times$ & 4.14 & 1.72 & -58.6\% \\
 & 1.0\% & 159 & 0.0647 & 0.0582 & 1.11$\times$ & 45.90 & 22.95 & -50.0\% \\
 & 0.5\% & 580 & 0.1812 & 0.1581 & 1.15$\times$ & 409.53 & 366.78 & -10.4\% \\
 & 0.3\% & 1,393 & 0.5690 & 0.4033 & \textbf{1.41$\times$} & 2125.28 & 506.08 & \textbf{-76.2\%} \\
\midrule
\multirow{6}{*}{\texttt{T10I4D100K} (Sparse)}
 & 4.0\% & 26 & 0.0399 & 0.0358 & 1.12$\times$ & 5.48 & 1.26 & -77.0\% \\
 & 3.0\% & 60 & 0.0653 & 0.0547 & 1.19$\times$ & 81.39 & 15.45 & -81.0\% \\
 & 2.0\% & 155 & 0.4125 & 0.1377 & 3.00$\times$ & 924.84 & 312.45 & -66.2\% \\
 & 1.5\% & 237 & 0.7651 & 0.2346 & \textbf{3.26$\times$} & 1719.26 & 1246.12 & -27.5\% \\
 & 1.0\% & 385 & 1.0615 & 0.5646 & 1.88$\times$ & 3751.47 & 1131.74 & -69.8\% \\
 & 0.7\% & 603 & 1.5323 & 0.8380 & 1.83$\times$ & 5486.69 & 1893.93 & -65.5\% \\
\bottomrule
\end{tabular}
\end{table*}

\subsection{Thực nghiệm 3: Số lượng Tập Phổ Biến (\textit{Output Size})}
Hình~\ref{fig:output_size} so sánh quy mô đầu ra trên thang logarit giữa hai nhóm dữ liệu. Ở nhóm \textbf{Dense} (\texttt{Chess}, \texttt{Mushroom}, \texttt{Accidents}), các item đồng xuất hiện với mật độ cao nên số tập phổ biến tăng theo hàm mũ khi giảm \texttt{min\_sup} (đường cong gần tuyến tính trên trục log: \texttt{Chess} tăng gần 180 lần từ 622 lên 111,239 tập; \texttt{Mushroom} tăng gần 200 lần từ 505 lên 99,071 tập). Ngược lại, ở nhóm \textbf{Sparse} (\texttt{Retail}, \texttt{T10I4D100K}), dù giảm \texttt{min\_sup} xuống $0.3\%$--$0.7\%$, số tập phổ biến vẫn chỉ ở mức vài trăm đến 1,393 tập. Hiện tượng này lý giải vì sao kỹ thuật \textit{Promotion} (Property 9) mang lại lợi ích đặc biệt lớn trên dữ liệu dày đặc, trong khi trên dữ liệu thưa phần lớn thời gian nằm ở pha xây \texttt{BMC-tree} và lọc các ứng viên dài 2--3.

\begin{figure}[htbp]
\centering
\includegraphics[width=\columnwidth]{figures/fig2_output_size.png}
\caption{Thực nghiệm 3: Số lượng tập phổ biến theo \texttt{min\_sup} (thang đo logarit) giữa dữ liệu Dày đặc (trái) và Thưa (phải).}
\label{fig:output_size}
\end{figure}

\subsection{Thực nghiệm 4: Mức sử dụng Bộ nhớ (\textit{Memory Consumption})}
Bảng~\ref{tab:exp4} và Hình~\ref{fig:memory} trình bày kết quả đo bộ nhớ ở mốc \texttt{min\_sup} trung bình cho từng tập dữ liệu. Bản \texttt{negfin\_optimized} giúp giảm tổng Heap cấp phát từ \textbf{27.52\% đến 69.55\%} trên cả 5 bộ dữ liệu (\texttt{Accidents} giảm 69.55\%, \texttt{Mushroom} giảm 60.00\%, \texttt{Retail} giảm 50.01\%). Đặc biệt, trên \texttt{T10I4D100K}, bộ nhớ vật lý đỉnh (\textit{Peak RSS}) giảm mạnh từ \textbf{449.08 MB xuống 372.50 MB} (tiết kiệm 76.58 MB RAM thực) nhờ biểu diễn cây LCRS giảm cấp phát cho nút lá và bố cục SoA gọn nhẹ hơn đáng kể so với mảng đối tượng \texttt{BitVector}.

\begin{table}[htbp]
\caption{Kết quả Thực nghiệm 4: So sánh Heap Allocation và Peak RSS giữa Baseline và Optimized}
\label{tab:exp4}
\centering
\resizebox{\columnwidth}{!}{%
\begin{tabular}{@{}lcrrrrr@{}}
\toprule
\textbf{Dataset} & \textbf{min\_sup} & \textbf{Alloc B} & \textbf{Alloc O} & \textbf{Giảm} & \textbf{RSS B} & \textbf{RSS O} \\
 & & \textbf{(MB)} & \textbf{(MB)} & \textbf{(\%)} & \textbf{(MB)} & \textbf{(MB)} \\
\midrule
\texttt{Chess} & 0.750 & 11.75 & 7.42 & -36.87 & 309.45 & 308.23 \\
\texttt{Mushroom} & 0.250 & 4.70 & 1.88 & -60.00 & 309.45 & 309.06 \\
\texttt{Accidents} & 0.600 & 25.61 & 7.80 & \textbf{-69.55} & 451.05 & 444.02 \\
\texttt{Retail} & 0.010 & 45.90 & 22.95 & -50.01 & 347.27 & 337.08 \\
\texttt{T10I4D100K} & 0.015 & 1719.26 & 1246.12 & -27.52 & 449.08 & \textbf{372.50} \\
\bottomrule
\end{tabular}%
}
\end{table}

\begin{figure}[htbp]
\centering
\includegraphics[width=\columnwidth]{figures/fig3_memory_comparison.png}
\caption{Thực nghiệm 4: So sánh tổng Heap cấp phát (trái) và RAM đỉnh Peak RSS (phải) giữa Baseline và Optimized.}
\label{fig:memory}
\end{figure}

\subsection{Thực nghiệm 5: Khả năng Mở rộng (\textit{Scalability})}
Bảng~\ref{tab:exp5} và Hình~\ref{fig:scalability} đánh giá khả năng mở rộng khi tăng kích thước CSDL theo các lát cắt $10\%, 25\%, 50\%, 75\%, 100\%$ trên \texttt{Retail} ($\text{min\_sup}=0.5\%$) và \texttt{Accidents} ($\text{min\_sup}=55\%$). Khi số lượng giao dịch tăng 10 lần, thời gian chạy bản tối ưu tăng \textbf{10.39$\times$} trên \texttt{Retail} (0.0146s $\rightarrow$ 0.1511s) và \textbf{10.60$\times$} trên \texttt{Accidents} (0.0352s $\rightarrow$ 0.3734s). Độ tăng trưởng gần như \textbf{tuyến tính $O(|DB|)$} xác nhận tính ổn định cao của cấu trúc \texttt{BMC-tree} và \texttt{NegNodeset} khi quy mô giao dịch tăng mạnh trong khi phân bố item giữ tương đối ổn định.

\begin{table}[htbp]
\caption{Kết quả Thực nghiệm 5: Khả năng mở rộng (10\% $\rightarrow$ 100\%) trên Retail và Accidents}
\label{tab:exp5}
\centering
\setlength{\tabcolsep}{4pt}
\begin{tabular}{@{}lcrrrr@{}}
\toprule
\textbf{Dataset} & \textbf{Tỷ lệ} & \textbf{\#Tx} & \textbf{\#Itemsets} & \textbf{Base (s)} & \textbf{Opt (s)} \\
\midrule
\multirow{5}{*}{\texttt{Retail} (0.5\%)}
 & 10\% & 8,816 & 692 & 0.0283 & 0.0146 \\
 & 25\% & 22,040 & 659 & 0.0471 & 0.0392 \\
 & 50\% & 44,081 & 613 & 0.1087 & 0.0726 \\
 & 75\% & 66,122 & 544 & 0.1395 & 0.1069 \\
 & 100\% & 88,162 & 580 & 0.1824 & 0.1511 \\
\midrule
\multirow{5}{*}{\texttt{Accidents} (55\%)}
 & 10\% & 34,018 & 3,628 & 0.0472 & 0.0352 \\
 & 25\% & 85,046 & 3,991 & 0.1109 & 0.0880 \\
 & 50\% & 170,092 & 4,160 & 0.2216 & 0.1847 \\
 & 75\% & 255,137 & 4,010 & 0.4289 & 0.2797 \\
 & 100\% & 340,183 & 4,051 & 0.4302 & 0.3734 \\
\bottomrule
\end{tabular}
\end{table}

\begin{figure}[htbp]
\centering
\includegraphics[width=\columnwidth]{figures/fig4_scalability.png}
\caption{Thực nghiệm 5: Thời gian chạy theo tỷ lệ kích thước CSDL trên \texttt{Retail} (trái) và \texttt{Accidents} (phải).}
\label{fig:scalability}
\end{figure}

\subsection{Thực nghiệm 6: Ảnh hưởng của Độ dài Giao dịch Trung bình}
Để kiểm chứng mối liên hệ giữa chiều dài giao dịch và cấu trúc \texttt{BMC-tree} trong lý thuyết \texttt{negFIN}, chúng tôi sinh 6 CSDL tổng hợp ($n_t = 10{,}000$, $|I| = 80$, $\text{min\_sup}=20\%$, phân bố lệch, \texttt{seed} cố định) với độ dài giao dịch trung bình tăng từ $10$ lên $30$. Kết quả tại Bảng~\ref{tab:exp6} và Hình~\ref{fig:txlen} cho thấy:
\begin{itemize}
    \item Khi \texttt{avg\_len} tăng từ $10 \rightarrow 30$, xác suất trùng tiền tố giữa các giao dịch giảm mạnh làm số nút \texttt{BMC-tree} phình to hơn \textbf{3,500 lần} (từ $63$ lên $222{,}694$ nút), kéo theo kích thước \texttt{NegNodeset} ($n$) và số tập phổ biến ($6 \rightarrow 833$) tăng phi tuyến.
    \item Chính ở kịch bản giao dịch dài (\texttt{avg\_len} $= 30$), bản tối ưu thể hiện ưu thế áp đảo: chỉ mất \textbf{0.0436s} so với \textbf{0.2128s} của bản cơ bản (\textbf{tăng tốc 4.88$\times$}), vì toàn bộ $n_f = 80$ bit được nén gọn trong một giá trị \texttt{UInt128} không cấp phát heap.
\end{itemize}

\begin{table}[htbp]
\caption{Kết quả Thực nghiệm 6: Ảnh hưởng của độ dài giao dịch trung bình ($10{,}000$ tx, 80 items, $\text{min\_sup}=20\%$)}
\label{tab:exp6}
\centering
\resizebox{\columnwidth}{!}{%
\begin{tabular}{@{}rrrrrrr@{}}
\toprule
\textbf{Dài TB} & \textbf{$n_f$} & \textbf{\#BMC Nodes} & \textbf{\#Itemsets} & \textbf{Base (s)} & \textbf{Opt (s)} & \textbf{Speedup} \\
\midrule
10 & 6 & 63 & 6 & 0.0033 & 0.0031 & 1.09$\times$ \\
14 & 16 & 7,393 & 19 & 0.0064 & 0.0051 & 1.25$\times$ \\
18 & 37 & 56,883 & 54 & 0.0233 & 0.0092 & 2.54$\times$ \\
22 & 69 & 140,668 & 134 & 0.0935 & 0.0164 & 5.72$\times$ \\
26 & 80 & 192,864 & 308 & 0.1321 & 0.0292 & 4.53$\times$ \\
30 & 80 & 222,694 & 833 & 0.2128 & 0.0436 & \textbf{4.88$\times$} \\
\bottomrule
\end{tabular}%
}
\end{table}

\begin{figure}[htbp]
\centering
\includegraphics[width=\columnwidth]{figures/fig5_tx_length.png}
\caption{Thực nghiệm 6: Ảnh hưởng độ dài giao dịch trung bình đến thời gian chạy (trái) và sự phình to của cây \texttt{BMC-tree} cùng số tập phổ biến (phải).}
\label{fig:txlen}
\end{figure}

\subsection{So sánh với SPMF (Java) và Đề xuất Hướng Tối ưu Tiếp theo}
So với cài đặt chuẩn của thư viện SPMF (Java), cài đặt \texttt{negFIN} bằng Julia có lợi thế biên dịch trực tiếp ra mã máy LLVM không bị overhead đóng gói đối tượng (\textit{object boxing}) của JVM, tận dụng trực tiếp mảng phẳng SoA và kiểu bit nguyên thủy trên thanh ghi. Ngược lại, hạn chế cố hữu của \texttt{negFIN} (như nhận định ở Mục 5.3.2 bài báo gốc~\cite{b9}) là trên dữ liệu rất thưa có $n_f$ lớn (\texttt{Retail} ở $\text{min\_sup} = 0.3\%$), mỗi nút \texttt{BMC-tree} mang một mã bitmap đầy đủ độ dài $\lceil n_f / 64 \rceil$ từ \texttt{UInt64}, tiêu tốn RAM nhiều hơn cây chỉ lưu thứ tự duyệt.

Từ các phân tích thực nghiệm trên, chúng tôi đề xuất \textbf{3 hướng tối ưu hóa tiếp theo}:
\begin{enumerate}
    \item \textbf{Song song hóa đa luồng (\texttt{Threads.@threads} / \texttt{@spawn}) ở tầng 1:} Các cây con xuất phát từ từng item $p \in 1..n_f$ ở tầng 1 chỉ \textit{đọc} mảng \texttt{nodesets} và độc lập hoàn toàn về trạng thái ghi cục bộ (\texttt{prefix}, \texttt{equiv\_stack}). Áp dụng lập lịch động giúp cân bằng tải tốt giữa các nhánh nặng (đầu $L_1$) và nhánh nhẹ (cuối $L_1$).
    \item \textbf{Nén Bitmap theo đoạn (\textit{Trimmed / Sparse Bitset}) cho dữ liệu thưa:} Theo \textit{Property 1}, toàn bộ \textit{don't-care section} luôn bằng $0$, và trên dữ liệu thưa phần lớn các word \texttt{UInt64} đầu cũng bằng $0$. Chỉ lưu dải word khác $0$ dưới dạng \texttt{(offset::UInt16, words::Vector\{UInt64\})} sẽ giảm mạnh Peak RAM trên \texttt{Retail} và \texttt{T10I4D100K}.
    \item \textbf{Khai thác Tập Đóng (\textit{Closed FIM}) và Tập Cực Đại (\textit{Maximal FIM}) trên \texttt{NegNodeset}:} Theo định hướng Mục 7 của Aryabarzan et al.~\cite{b9}, tích hợp kiểm tra tính đóng trực tiếp trên tập \texttt{equivalent\_items} giúp tránh bùng nổ đầu ra trên các bộ dữ liệu dày đặc ở ngưỡng \texttt{min\_sup} rất thấp.
\end{enumerate}

% ==============================================================================
\section{Chương 3: Ứng Dụng Thực Tế Với Mã Nguồn \texttt{negFIN} Thuần Julia}
% ==============================================================================

\subsection{Kiến trúc Tích hợp Thuần Julia}
Nguyên tắc cốt lõi của Chương 3 là \textbf{toàn bộ tập phổ biến ở cả hai ứng dụng đều được khai thác trực tiếp bằng hàm \texttt{negfin\_optimized} (và đối chiếu được với \texttt{negfin\_baseline}) tự viết của nhóm} trong \texttt{src/algorithm/}, tuyệt đối không sử dụng thư viện FIM ngoài hay mã ứng dụng Python. Module \texttt{src/applications/applications.jl} nạp trực tiếp mã nguồn Chương 1 và cung cấp hai ứng dụng độc lập cùng bộ kiểm thử tự động \texttt{tests/test\_ch3\_julia.jl} (\textbf{27/27 PASS}). Mọi kết quả đều trải qua bước kiểm chứng chéo \texttt{validate\_support} đếm trực tiếp trên giao dịch gốc (\texttt{Support cross-check: PASS}). Hai tệp dữ liệu nguyên bản (\texttt{groceries.csv} từ gói R \texttt{arules}~\cite{b10} và \texttt{digits8x8.csv} từ UCI ML Repository) cũng được lưu trữ dưới dạng bản sao lưu công khai (\textit{public mirror}) tại cùng kho Hugging Face của nhóm để tự động tải khi chạy trên Google Colab.

\subsection{Ứng dụng 1: Phân tích Giỏ hàng (\textit{Market Basket Analysis})}

\subsubsection{Dữ liệu và Quy trình Xử lý}
Chúng tôi sử dụng tập dữ liệu bán lẻ thực tế \textbf{\texttt{Groceries}}~\cite{b10} gồm \textbf{9,835 giao dịch} thực tế tại một cửa hàng tạp hóa với \textbf{169 nhóm sản phẩm} phân biệt. Quy trình xử lý trong \texttt{src/applications/market\_basket.jl} gồm 4 bước:
\begin{enumerate}
    \item \textbf{Chuẩn hóa và mã hóa:} Loại bỏ ô trống, chuẩn hóa khoảng trắng, loại item trùng trong cùng hóa đơn và ánh xạ 169 sản phẩm theo thứ tự chữ cái sang mã nguyên $1..169$ (lưu tại \texttt{item\_mapping.csv} và \texttt{transactions.spmf}).
    \item \textbf{Khai thác tập phổ biến bằng \texttt{negFIN}:} Chạy \texttt{negfin\_optimized} với $\text{min\_sup} = 0.01$ ($\ge 99$ hóa đơn), thu được \textbf{333 tập phổ biến} trong \textbf{0.2314 giây} (\texttt{Support cross-check: PASS}).
    \item \textbf{Tự cài đặt sinh luật kết hợp ($X \rightarrow Y$):} Với mỗi tập phổ biến $Z$ có $|Z| \ge 2$, duyệt mọi phân hoạch không rỗng $X \subset Z$, $Y = Z \setminus X$ và tính các độ đo:
    \begin{align}
    \text{conf}(X \!\rightarrow\! Y) &= \frac{\text{sup}(X \cup Y)}{\text{sup}(X)}, \label{eq:conf} \\
    \text{lift}(X \!\rightarrow\! Y) &= \frac{\text{conf}(X \!\rightarrow\! Y)}{\text{sup}(Y)} = \frac{\text{sup}(X \cup Y)}{\text{sup}(X)\text{sup}(Y)}, \label{eq:lift} \\
    \text{lev}(X \!\rightarrow\! Y) &= \text{sup}(X \cup Y) - \text{sup}(X)\text{sup}(Y). \label{eq:lev}
    \end{align}
    \item \textbf{Lọc và xếp hạng:} Với ngưỡng $\text{min\_conf} = 0.15$, hệ thống sinh được \textbf{322 luật kết hợp} hợp lệ và xếp hạng Top-10 theo chỉ số \texttt{Lift} giảm dần.
\end{enumerate}

\subsubsection{Kết quả Top-10 Luật theo Lift và Diễn giải Kinh doanh}
Bảng~\ref{tab:top10_lift} và Hình~\ref{fig:ch3_lift} trình bày đầy đủ Top-10 luật kết hợp có chỉ số \texttt{Lift} cao nhất.

\begin{table*}[t]
\caption{Ứng dụng 1: Top-10 luật kết hợp theo chỉ số Lift trên tập dữ liệu Groceries (9,835 giao dịch, $\text{min\_sup}=0.01$, $\text{min\_conf}=0.15$)}
\label{tab:top10_lift}
\centering
\begin{tabular}{@{}cllrrrrr@{}}
\toprule
\textbf{Hạng} & \textbf{Vế trái (Antecedent $X$)} & \textbf{Vế phải (Consequent $Y$)} & \textbf{Count} & \textbf{Support} & \textbf{Confidence} & \textbf{Lift} & \textbf{Leverage} \\
\midrule
1 & \texttt{curd} & \texttt{whole milk + yogurt} & 99 & 0.0101 & 18.89\% & \textbf{3.37} & 0.0071 \\
2 & \texttt{whole milk + yogurt} & \texttt{curd} & 99 & 0.0101 & 17.97\% & \textbf{3.37} & 0.0071 \\
3 & \texttt{citrus fruit + other vegetables} & \texttt{root vegetables} & 102 & 0.0104 & 35.92\% & \textbf{3.30} & 0.0072 \\
4 & \texttt{other vegetables + yogurt} & \texttt{whipped/sour cream} & 100 & 0.0102 & 23.42\% & \textbf{3.27} & 0.0071 \\
5 & \texttt{other vegetables + tropical fruit} & \texttt{root vegetables} & 121 & 0.0123 & 34.28\% & \textbf{3.14} & 0.0084 \\
6 & \texttt{beef} & \texttt{root vegetables} & 171 & 0.0174 & 33.14\% & \textbf{3.04} & 0.0117 \\
7 & \texttt{root vegetables} & \texttt{beef} & 171 & 0.0174 & 15.95\% & \textbf{3.04} & 0.0117 \\
8 & \texttt{citrus fruit + root vegetables} & \texttt{other vegetables} & 102 & 0.0104 & \textbf{58.62\%} & \textbf{3.03} & 0.0069 \\
9 & \texttt{root vegetables + tropical fruit} & \texttt{other vegetables} & 121 & 0.0123 & \textbf{58.45\%} & \textbf{3.02} & 0.0082 \\
10 & \texttt{other vegetables + whole milk} & \texttt{root vegetables} & 228 & 0.0232 & 30.98\% & \textbf{2.84} & \textbf{0.0150} \\
\bottomrule
\end{tabular}
\end{table*}

\begin{figure}[htbp]
\centering
\includegraphics[width=\columnwidth]{figures/ch3_top10_lift.png}
\caption{Ứng dụng 1: Biểu đồ Top-10 luật kết hợp theo chỉ số Lift trên bộ dữ liệu \texttt{Groceries}.}
\label{fig:ch3_lift}
\end{figure}

Từ kết quả thực nghiệm, ta rút ra 3 cụm hành vi mua sắm nổi bật có giá trị ứng dụng kinh doanh rõ rệt:
\begin{itemize}
    \item \textbf{Cụm chế phẩm từ sữa (Hạng 1, 2, 4):} Luật \texttt{\{curd\} $\Rightarrow$ \{whole milk, yogurt\}} đạt $\text{Lift} = 3.37$, cho thấy khách hàng mua sữa đông có xác suất mua kèm combo sữa tươi và sữa chua cao gấp 3.37 lần mức ngẫu nhiên. Tương tự, \texttt{\{other vegetables, yogurt\} $\Rightarrow$ \{whipped/sour cream\}} đạt $\text{Lift} = 3.27$. \textit{Gợi ý kinh doanh:} Bố trí liền kề quầy tủ mát sữa đông, sữa chua, kem chua và thiết kế gói ưu đãi bữa sáng.
    \item \textbf{Cụm rau củ quả tươi (Hạng 3, 5, 8, 9, 10):} Luật \texttt{\{citrus fruit, root vegetables\} $\Rightarrow$ \{other vegetables\}} và \texttt{\{root vegetables, tropical fruit\} $\Rightarrow$ \{other vegetables\}} có độ tin cậy rất cao (\textbf{58.62\%} và \textbf{58.45\%}) với $\text{Lift} > 3.0$; đặc biệt luật hạng 10 (\texttt{\{other vegetables, whole milk\} $\Rightarrow$ \{root vegetables\}}) đạt độ phủ lớn nhất trong top-10 (228 hóa đơn, $\text{Leverage} = 0.0150$). \textit{Gợi ý:} Sắp xếp khu trái cây họ cam/nhiệt đới dẫn trực tiếp sang khu rau củ để tối ưu hóa hành trình mua thực phẩm tươi sống hàng ngày.
    \item \textbf{Cụm nguyên liệu nấu ăn (Hạng 6, 7):} Luật \texttt{\{beef\} $\Rightarrow$ \{root vegetables\}} (171 hóa đơn, $\text{Confidence} = 33.14\%$, $\text{Lift} = 3.04$) phản ánh thói quen mua thịt bò kèm rau củ hầm/nấu súp. \textit{Gợi ý:} Triển khai trưng bày chéo (\textit{cross-merchandising}) đặt khay củ quả sơ chế cạnh quầy thịt bò.
\end{itemize}
\textit{Lưu ý khoa học:} Chỉ số $\text{Lift} > 1$ chỉ phản ánh \textbf{mối tương quan thống kê dương} trên dữ liệu lịch sử, không tự động chứng minh quan hệ nhân quả hay đảm bảo tăng lợi nhuận; mọi đề xuất trưng bày hoặc khuyến mãi cần được kiểm chứng thực tế bằng thử nghiệm A/B.

\subsection{Ứng dụng 2: Khai thác Mẫu Trực quan trên Hình ảnh (\textit{Bag of Visual Words})}

\subsubsection{Dữ liệu và Quy trình Chuyển đổi Ảnh thành Giao dịch}
Lấy cảm hứng từ mô hình \textit{Bag of Visual Words} (BoVW) của Sivic \& Zisserman~\cite{b11}, ứng dụng thứ hai (\texttt{src/applications/visual\_words.jl}) áp dụng \texttt{negFIN} trên tập ảnh chữ số viết tay \textbf{\texttt{UCI Optical Recognition of Handwritten Digits (8$\times$8)}} (4 lớp chữ số \texttt{0, 1, 3, 8}, tối đa 150 ảnh/lớp, tổng cộng 600 ảnh -- Hình~\ref{fig:ch3_digits}). Quy trình hoàn toàn thuần Julia và tránh rò rỉ dữ liệu (\textit{data leakage}):
\begin{enumerate}
    \item \textbf{Chia tập dữ liệu:} Phân tầng theo từng lớp tỷ lệ 75\%/25\% với \texttt{seed = 42}, thu được \textbf{448 ảnh huấn luyện} và \textbf{152 ảnh kiểm thử}. Nhãn lớp \textbf{không} được sử dụng trong pha học từ điển và pha khai thác \texttt{negFIN}.
    \item \textbf{Trích xuất đặc trưng cục bộ:} Do ảnh có độ phân giải nhỏ $8 \times 8$, thay cho SIFT/ORB vốn cần ảnh lớn, chúng tôi tự cài đặt bộ trích \textbf{9 dense gradient descriptors} trên lưới $3 \times 3$ vùng cục bộ cho mỗi ảnh. Mỗi descriptor là vector 6 chiều gồm: cường độ trung bình, độ lệch chuẩn và 4 thành phần độ lớn gradient có hướng ($+x, -x, +y, -y$).
    \item \textbf{Học từ điển \textit{Codebook} bằng K-means tự cài:} Thuật toán K-means thuần Julia ($K = 16$, \texttt{seed = 42}) được huấn luyện chỉ trên các descriptor của 448 ảnh train, sinh ra 16 \textit{visual words} ($W_1 \dots W_{16}$). Mỗi ảnh được lượng tử hóa thành một giao dịch chứa tập các visual word ID xuất hiện trong ảnh (đã loại trùng lặp).
    \item \textbf{Khai thác mẫu trực quan bằng \texttt{negFIN}:} Chạy \texttt{negfin\_optimized} trên 448 giao dịch train với $\text{min\_sup} = 0.12$ trong \textbf{0.2692 giây}, tìm được \textbf{154 tập phổ biến} (\texttt{Support cross-check: PASS}) và chọn lọc \textbf{60 mẫu trực quan} có độ dài $\ge 2$.
\end{enumerate}

\begin{figure}[htbp]
\centering
\includegraphics[width=0.92\columnwidth]{figures/ch3_digits_samples.png}
\caption{Ứng dụng 2: Một số mẫu ảnh chữ số viết tay $8 \times 8$ đại diện cho 4 lớp (\texttt{0, 1, 3, 8}) trong tập dữ liệu \texttt{UCI Digits}.}
\label{fig:ch3_digits}
\end{figure}

\subsubsection{Phân tích Mẫu Đặc trưng theo Lớp (\textit{Enrichment}) và Đánh giá Phân loại}
Sau khi khai thác không giám sát, nhãn tập train được dùng để tính độ làm giàu đặc trưng (\textit{enrichment}) của từng mẫu $P$ đối với lớp $C$:
\begin{equation}
\text{enrichment}(P, C) = \frac{P(P \subseteq T \mid \text{class}(T) = C)}{P(P \subseteq T)}.
\label{eq:enrichment}
\end{equation}
Với 4 lớp cân bằng, giá trị lý thuyết tối đa của $\text{enrichment}$ là $4.00$ (mẫu chỉ xuất hiện độc quyền trong đúng 1 lớp). Bảng~\ref{tab:ch3_patterns} và Hình~\ref{fig:ch3_enriched} cho thấy \texttt{negFIN} phát hiện được nhiều tổ hợp từ trực quan có tính phân biệt lớp cực cao: chẳng hạn tổ hợp \texttt{W6+W7+W10} và \texttt{W2+W6+W7+W10} xuất hiện trong gần 60\% ảnh chữ số \texttt{0} với $\text{enrichment} = \mathbf{4.00}$ (gần như độc quyền cho lớp \texttt{0}), hay tổ hợp \texttt{W3+W4+W13+W16} đặc trưng mạnh cho lớp \texttt{3} với $\text{enrichment} = \mathbf{3.64}$.

\begin{table}[htbp]
\caption{Ứng dụng 2: Một số mẫu trực quan có độ làm giàu (Enrichment) cao nhất trên tập huấn luyện}
\label{tab:ch3_patterns}
\centering
\resizebox{\columnwidth}{!}{%
\begin{tabular}{@{}lcrrrr@{}}
\toprule
\textbf{Mẫu Visual Words} & \textbf{Lớp} & \textbf{Count} & \textbf{Global Sup} & \textbf{Class Sup} & \textbf{Enrich.} \\
\midrule
\texttt{W6+W7+W10} & 0 & 66 & 0.1473 & 0.5893 & \textbf{4.00} \\
\texttt{W2+W6+W7+W10} & 0 & 65 & 0.1451 & 0.5804 & \textbf{4.00} \\
\texttt{W2+W6+W10} & 0 & 97 & 0.2188 & 0.8661 & 3.96 \\
\texttt{W6+W10} & 0 & 100 & 0.2321 & 0.8929 & 3.85 \\
\texttt{W3+W4+W13+W16} & 3 & 51 & 0.1250 & 0.4554 & \textbf{3.64} \\
\texttt{W3+W4+W16} & 3 & 53 & 0.1317 & 0.4732 & 3.59 \\
\bottomrule
\end{tabular}%
}
\end{table}

\begin{figure}[htbp]
\centering
\includegraphics[width=\columnwidth]{figures/ch3_enriched_patterns.png}
\caption{Ứng dụng 2: Các mẫu trực quan phổ biến có độ làm giàu (\textit{enrichment}) cao nhất phân theo lớp chữ số.}
\label{fig:ch3_enriched}
\end{figure}

Để kiểm tra tác dụng thực tế trên bài toán phân loại ảnh, chúng tôi đánh giá bộ phân loại \textit{Nearest-Centroid} tự cài đặt trên tập kiểm thử độc lập (152 ảnh) với hai biểu diễn đặc trưng (Bảng~\ref{tab:ch3_acc}): (i) chỉ gồm 16 tần suất visual words (độ chính xác \textbf{79.61\%}), và (ii) ghép thêm 60 đặc trưng nhị phân chỉ báo sự xuất hiện của các frequent patterns (tổng 76 chiều, độ chính xác \textbf{78.95\%}).

\begin{table}[htbp]
\caption{So sánh độ chính xác phân loại Nearest-Centroid trên tập kiểm thử độc lập (152 ảnh)}
\label{tab:ch3_acc}
\centering
\begin{tabular}{@{}lcc@{}}
\toprule
\textbf{Biểu diễn đặc trưng} & \textbf{Số chiều} & \textbf{Test Accuracy} \\
\midrule
Visual Words đơn lẻ (BoVW cơ bản) & 16 & \textbf{79.61\%} (0.7961) \\
Visual Words + 60 Frequent Patterns & 76 & 78.95\% (0.7895) \\
\bottomrule
\end{tabular}
\end{table}

\textbf{Phân tích khoa học và Giới hạn phương pháp:} Việc bổ sung trực tiếp toàn bộ 60 mẫu phổ biến làm giảm nhẹ độ chính xác khoảng 0.66\% (tương đương 1 ảnh trên 152 ảnh test). Đây là hiện tượng hoàn toàn hợp lý về mặt lý thuyết học máy: (1) nhiều tập phổ biến lồng nhau (ví dụ \texttt{W6+W10}, \texttt{W2+W6+W10}, \texttt{W2+W6+W7+W10}) có độ tương quan rất cao, gây lặp trọng số trong khoảng cách Euclidean không trọng số của bộ phân loại \textit{Nearest-Centroid}; và (2) quá trình khai thác \texttt{negFIN} là không giám sát nên không tối ưu trực tiếp cho biên phân loại như các phương pháp chọn đặc trưng phân biệt (\textit{discriminative pattern selection}). Ngoài ra, ảnh \texttt{Digits 8$\times$8} có kích thước rất nhỏ và descriptor gradient tự cài mang tính minh họa quy trình, khác biệt đáng kể so với bài toán ảnh tự nhiên độ phân giải cao (như Caltech-101 hay CIFAR). Để cải thiện độ chính xác phân loại thực tế, cần kết hợp lọc mẫu đóng (\textit{Closed Patterns}) theo tiêu chí \textit{information gain} và sử dụng bộ phân loại có điều chuẩn như SVM hoặc Logistic Regression.

\subsection{Khả năng Tái lập (Reproducibility) và Tài nguyên Công khai}
Toàn bộ mã nguồn, bộ kiểm thử, số liệu thực nghiệm và notebook demo trên Google Colab được công khai để đảm bảo khả năng tái lập 100\% theo đúng quy chuẩn đồ án (Bảng~\ref{tab:repro}).

\begin{table}[htbp]
\caption{Danh mục tài nguyên công khai phục vụ tái thực hiện đồ án}
\label{tab:repro}
\centering
\resizebox{\columnwidth}{!}{%
\begin{tabular}{@{}ll@{}}
\toprule
\textbf{Thành phần} & \textbf{Đường dẫn / Lệnh thực thi} \\
\midrule
Mã nguồn GitHub & \texttt{github.com/Thanhtung46/Frequent\_Itemset\_Mining} \\
Dataset Mirror & \texttt{huggingface.co/datasets/stephent9/} \\
 & \texttt{Lab1\_Frequent\_Itemset\_Mining} \\
Unit Tests & \texttt{julia --project=. tests/runtests.jl} \\
 & \texttt{julia --project=. tests/test\_ch3\_julia.jl} \\
CLI chuẩn SPMF & \texttt{julia --project=. src/main.jl --input ...} \\
Colab Ch. 2 & \texttt{notebooks/colab/ch2\_experiments.ipynb} \\
Colab Ch. 3 & \texttt{notebooks/colab/04\_ch3\_all\_applications.ipynb} \\
Local NBs & \texttt{notebooks/local/*.ipynb} (Pure Julia) \\
\bottomrule
\end{tabular}%
}
\end{table}

% ==============================================================================
\section{Kết Luận}
% ==============================================================================
Đồ án đã hoàn thành toàn diện cả ba chương nghiên cứu, cài đặt và ứng dụng thuật toán \texttt{negFIN} bằng ngôn ngữ \textbf{Julia}:
\begin{enumerate}
    \item Cài đặt chính xác 100\% thuật toán \texttt{negFIN} theo bài báo gốc (Aryabarzan et al., 2018) với đầy đủ 4 bậc, đạt 56/56 kiểm thử tự động thành công, kèm phiên bản tối ưu hóa Julia-native tận dụng Multiple Dispatch bit-packing, LCRS \texttt{BMC-tree}, SoA \texttt{NodeSet}, Early-Abort và xuất Promotion không cấp phát \texttt{FIS\_parent}.
    \item Hoàn tất 6 thực nghiệm định lượng trên 5 bộ dữ liệu chuẩn công khai tại Hugging Face và dữ liệu tổng hợp, chứng minh tốc độ tăng lên đến \textbf{3.26$\times$} trên benchmark thực (\textbf{4.88$\times$} trên giao dịch dài), giảm \textbf{27.5\%--69.5\%} Heap cấp phát và mở rộng tuyến tính $O(|DB|)$.
    \item Triển khai thành công hai ứng dụng thực tế thuần Julia gọi trực tiếp mã nguồn \texttt{negFIN} của nhóm trên dữ liệu \texttt{Groceries} (9,835 giao dịch, Top-10 Lift đạt 3.37) và ảnh \texttt{UCI Digits} (BoVW + K-means tự cài, Enrichment đạt 4.00), có đầy đủ notebook Google Colab tái lập 100\%.
\end{enumerate}

Trong tương lai, nhóm dự kiến triển khai song song hóa đa luồng \texttt{Threads.@threads} ở tầng 1, nén \textit{Trimmed Bitset} cho dữ liệu thưa quy mô hàng chục nghìn item, và mở rộng sang khai thác tập phổ biến đóng (\textit{Closed FIM}) để phục vụ tốt hơn cho bài toán phân loại dữ liệu chiều cao.

% ==============================================================================
\begin{thebibliography}{00}

\bibitem{b1} R. Agrawal, T. Imieliński, and A. Swami, ``Mining association rules between sets of items in large databases,'' \textit{ACM SIGMOD Record}, vol. 22, no. 2, pp. 207--216, Jun. 1993.

\bibitem{b2} R. Agrawal and R. Srikant, ``Fast algorithms for mining association rules in large databases,'' in \textit{Proc. 20th Int. Conf. Very Large Data Bases (VLDB)}, 1994, pp. 487--499.

\bibitem{b3} J. Han, J. Pei, and Y. Yin, ``Mining frequent patterns without candidate generation,'' \textit{ACM SIGMOD Record}, vol. 29, no. 2, pp. 1--12, Jun. 2000.

\bibitem{b4} M. J. Zaki, ``Scalable algorithms for association mining,'' \textit{IEEE Trans. Knowl. Data Eng.}, vol. 12, no. 3, pp. 372--390, May/Jun. 2000.

\bibitem{b5} Z.-H. Deng and Z. Wang, ``A new fast vertical method for mining frequent patterns,'' \textit{Int. J. Comput. Intell. Syst.}, vol. 3, no. 6, pp. 733--744, Dec. 2010.

\bibitem{b6} Z.-H. Deng, Z. Wang, and J. Jiang, ``A new algorithm for fast mining frequent itemsets using N-lists,'' \textit{Sci. China Inf. Sci.}, vol. 55, no. 9, pp. 2008--2030, Sep. 2012.

\bibitem{b7} Z.-H. Deng and S.-L. Lv, ``Fast mining frequent itemsets using nodesets,'' \textit{Expert Syst. Appl.}, vol. 41, no. 10, pp. 4505--4512, Aug. 2014.

\bibitem{b8} Z.-H. Deng, ``DiffNodesets: An efficient structure for fast mining frequent itemsets,'' \textit{Appl. Soft Comput.}, vol. 41, pp. 214--223, Apr. 2016.

\bibitem{b9} N. Aryabarzan, B. Minaei-Bidgoli, and M. Teshnehlab, ``negFIN: An efficient algorithm for fast mining frequent itemsets,'' \textit{Expert Syst. Appl.}, vol. 105, pp. 129--143, Sep. 2018.

\bibitem{b10} M. Hahsler, K. Hornik, and T. Reutterer, ``Implications of probabilistic data modeling for mining association rules,'' in \textit{From Data and Information Analysis to Knowledge Engineering}, Springer, 2006, pp. 598--605.

\bibitem{b11} J. Sivic and A. Zisserman, ``Video Google: A text retrieval approach to object matching in videos,'' in \textit{Proc. 9th IEEE Int. Conf. Computer Vision (ICCV)}, 2003, pp. 1470--1477.

\end{thebibliography}

\end{document}
